% =====================================================================
%  Guide — English sample (verifies all template features)
%  Compile with XeLaTeX:
%      xelatex main.tex   (run 2x for the TOC)
%
%  Class options:
%    portuguese  — Portuguese labels (default: English)
%    indentbody  — indent body text by \contentindent (off by default)
%    notime      — hide compilation time on cover page
%    nochangelog — disable changelog rendering
% =====================================================================
\documentclass{guide}

% --- Document configuration -----------------------------------------
\setguidetitle{Guide: Getting Started with Huawei Cloud ECS}
\setheadertitle{Huawei Cloud -- ECS Quick Start Guide}
\setcovertext{Huawei Technologies CO., LTD}
\setdocversion{1.6.0}
\setdocdate{\today}
\setheaderlogo{common-assets/huawei-logo-header.png}
\setcoverlogo{common-assets/huawei-logo-cover.png}

\begin{document}

% ---- COVER ---------------------------------------------------------
\makecover

% ---- TABLE OF CONTENTS ---------------------------------------------
\maketoc

% ---- BODY (restarts page numbering at 1) ---------------------------
\startbody

\begin{infobox}
This document demonstrates the \inlinecode{guide} template (\inlinecode{guide.cls}) for
Huawei Cloud guides. It exercises all available commands and environments:
cover page, table of contents, objectives, code blocks, images, callouts,
badge, changelog, and the \inlinecode{\textbackslash menu} command for menu paths.
\end{infobox}

% =====================================================================
%  Chapter 1 — Provision an Elastic Cloud Server
% =====================================================================
\section{Provision an Elastic Cloud Server}

\begin{objectives}
  \generalobjective{This lab walks you through provisioning an Elastic Cloud Server
    (ECS) on Huawei Cloud via the Console. By the end you will be able to
    launch a virtual machine, connect to it over SSH, and verify that it is
    running.}

  \prerequisites
  \begin{itemize}
    \item Active IAM account on Huawei Cloud (valid username and password).
    \item A web browser (Google Chrome, Mozilla Firefox, or Microsoft Edge).
    \item An SSH client installed locally (OpenSSH, PuTTY, etc.).
    \item A key pair created in the Console, or permission to create one.
  \end{itemize}
\end{objectives}

\begin{infobox}
This guide covers the basic ECS provisioning flow. For advanced operations,
consult the official Huawei Cloud documentation.
\end{infobox}

\subsection{Launch the instance}

\objective{Create a single ECS instance in the \inlinecode{sa-brazil-1} region
  using the Console wizard.}

\stepbystep
\begin{enumerate}
  \item In the Console, open the service catalog and navigate to
        \menu{Console, Elastic Cloud Server}.

  \imagecap[width=0.8\linewidth]{assets/exemplo-menu.png}{Service catalog in the Console.}

  \item Click \textbf{Create ECS} and fill in the basic configuration:

  \begin{code}
Name:           web-server-01
Flavor:         s6.small.1 (1 vCPU, 1 GB RAM)
Image:          Ubuntu 22.04 Server 64bit
Disk:           40 GB General Purpose SSD
VPC:            vpc-default
Subnet:         subnet-default
Security Group: sg-ssh (allows inbound TCP 22)
  \end{code}

  \item Under \textbf{Login Configuration}, select the key pair you created
        earlier. Click \textbf{Create Now} and wait for the status to change
        to \textbf{Running}.

  \image[width=0.5\linewidth]{assets/exemplo-login.png}

  \note{The instance is assigned a public IP (EIP). If you need a fixed IP,
    allocate a dedicated EIP and bind it after creation.}

  \begin{warning}
  The EIP assigned during creation is released when the instance is deleted.
  Allocate a dedicated EIP if you need to keep the address.
  \end{warning}
\end{enumerate}

\subsection{Connect over SSH}

\subsubsection{Retrieve the public IP}

On the ECS details page, copy the \textbf{EIP} (Elastic IP) field. This is
the address you will use to connect.

\paragraph{Using OpenSSH}
From a terminal, connect with the private key that matches the key pair
selected during creation:

\begin{code}[bash]
chmod 400 ~/.ssh/huawei_key.pem
ssh -i ~/.ssh/huawei_key.pem ubuntu@<EIP>
\end{code}

\paragraph{Using PuTTY}
Convert the key to PPK format with PuTTYgen, then open a session to the
EIP with the converted key.

\subsection{Verify the instance}

Once connected, run a few checks:

\begin{code}[bash]
uname -a          # confirm OS version
df -h             # check disk space
systemctl status  # verify services are running
\end{code}

Alternatively, save the checks to a script and run it with
\codefile[bash]{assets/example-script.sh}

\subsubsection{Recommended flavors}

The table below lists recommended ECS flavors for development workloads:

\begin{table}[H]
  \begin{hutable}{|l|l|l|}
    \rowcolor{huaweired} \thd{Flavor} & \thd{vCPU / RAM} & \thd{Use case} \\
    \tbody
    ac8.large.2  & 2 vCPU / 4 GB  & Light development \\
    ac8.xlarge.2 & 4 vCPU / 8 GB  & Recommended --- full coding stack \\
    ac8.2xlarge.2 & 8 vCPU / 16 GB & Heavy workloads + multiple agents \\
  \end{hutable}
  \caption{Recommended ECS flavors for development.}
\end{table}

\imageplaceholder{assets/ecs-flavors.png}{Screenshot of the ECS flavor selection page in the Console}

\note{This is an English sample that exercises every command and
  environment provided by the \inlinecode{guide} class. Use it as a reference
  when writing your own guides.}

% =====================================================================
%  Chapter 2 — Configure the Terraform Provider
% =====================================================================
\section{Configure the Terraform Provider}

\begin{objectives}
  \generalobjective{This chapter shows how to obtain an AK/SK pair from the Console
    and declare the Huawei Cloud provider in a Terraform configuration file,
    preparing the environment for Infrastructure as Code.}

  \prerequisites
  \begin{itemize}
    \item Chapter 1 completed (Console access working).
    \item IAM permission to create access credentials.
    \item \inlinecode{terraform} installed locally (version 1.0 or later).
  \end{itemize}
\end{objectives}

\subsection{Create an AK/SK}

\objective{Generate an Access Key ID / Secret Access Key pair to authenticate
  external tools with Huawei Cloud.}

In the Console, click your username in the top-right corner and open
\textbf{My Credentials}. In the \textbf{Access Keys} section, click
\textbf{Create Access Key} and download the generated file. Store it in a
safe location: the Secret Access Key is shown only once.

\subsection{Declare the provider}

\stepbystep
\begin{enumerate}
  \item Create a file named \param{provider.tf} in your project directory
        with the following content:

  \begin{code}
terraform {
  required_providers {
    huaweicloud = {
      source  = "huaweicloud/huaweicloud"
      version = ">= 1.60.0"
    }
  }
}

provider "huaweicloud" {
  region     = "sa-brazil-1"
  access_key = "YOUR_ACCESS_KEY_HERE"
  secret_key = "YOUR_SECRET_KEY_HERE"
}
  \end{code}

  Replace the \param{access\_key} and \param{secret\_key} values with the
  credentials obtained in the previous section.

  \item In a terminal, in the project folder, initialize the Terraform
        working directory:

  \begin{code}[bash]
terraform init
terraform plan
  \end{code}

  \note{Never commit the \param{provider.tf} file with real credentials.
    Use environment variables or a state file ignored by version control.}
\end{enumerate}

\subsection{Next steps}

With the provider declared, you can start defining resources in
\inlinecode{terraform}. For the full list of available resources, see the official
documentation at
\weblink{https://registry.terraform.io/providers/huaweicloud/huaweicloud/latest/docs}{Huawei Cloud Provider}.

\begin{tip}
Use \inlinecode{terraform validate} to check syntax before applying changes.
\end{tip}

\badge{Example}

% =====================================================================
%  Chapter 3 — Clean Up
% =====================================================================
\section{Clean Up}

\begin{objectives}
  \generalobjective{Remove the resources created in this lab to avoid ongoing charges.}
\end{objectives}

\stepbystep
\begin{enumerate}
  \item In the Console, navigate to \textbf{Elastic Cloud Server}.
  \item Select the instance \inlinecode{web-server-01}, click \textbf{More} and
        choose \textbf{Delete}.
  \item Confirm deletion. Optionally release the EIP if it was dedicated.
\end{enumerate}

\note{If you used Terraform, run \inlinecode{terraform destroy} instead to remove
  all resources declared in the configuration.}

\begin{changelog}
  \changelogentry{1.7.0}{\today}{\item Updated to base template v2.16.0 — Lua filter fixes and build system improvements.}
  \changelogentry{1.6.0}{2026-08-12}{
    \item The \texttt{changelog} environment now emits its own section heading; \texttt{[nochangelog]} suppresses the heading and entries in one switch.
  }
  \changelogentry{1.5.1}{2026-08-12}{
    \item Fix \texttt{hutable} centering and row-color fill: use \texttt{\textbackslash hline} so \texttt{colortbl} row colors fill cleanly up to the rules.
  }
  \changelogentry{1.5.0}{2026-08-12}{
    \item New \texttt{hutable} environment: full-grid tables with Huawei-red header and alternating body rows.
    \item Floats now default to \texttt{[H]} placement so tables and figures appear in source order.
  }
  \changelogentry{1.4.1}{2026-08-10}{
    \item H1 section title size reduced from 27pt to 20pt.
  }
  \changelogentry{1.4.0}{2026-08-10}{
    \item Default image size increased: width 65\% $\to$ 90\% of text width,
          height 40\% $\to$ 50\% of text height.
  }
  \changelogentry{1.3.0}{2026-08-10}{
    \item Enlarged H1 section title (18pt $\to$ 27pt, 50\% bigger).
    \item \texttt{nochangelog} option now hides version, date, and time on the
          cover page.
    \item Document date defaults to \inlinecode{\textbackslash today}
          (compilation date).
  }
  \changelogentry{1.2.1}{2026-08-09}{
    \item Rolled back \texttt{/ActualText} markers — broke code block rendering.
  }
  \changelogentry{1.2.0}{2026-08-09}{
    \item Added key-value sizing options to \inlinecode{\textbackslash image} and
          \inlinecode{\textbackslash imagecap} — \inlinecode{width} and
          \inlinecode{height} can now be set independently.
  }
  \changelogentry{1.1.3}{2026-08-09}{
    \item Fixed PDF copy-paste of URLs in code blocks — disabled Cascadia Code
          programming ligatures (\texttt{calt}) that corrupted \texttt{//} into
          \texttt{))} on copy.
  }
  \changelogentry{1.1.2}{2026-08-09}{
    \item Fixed italic font rendering — HarmonyOS Sans and Cascadia Code have
          no italic variants; enabled synthetic slant (\texttt{AutoFakeSlant}).
  }
  \changelogentry{1.1.1}{2026-08-08}{
    \item Changed table and figure caption labels to black (was Huawei red).
  }
  \changelogentry{1.1.0}{2026-08-08}{
    \item Added Huawei-branded table styling: red header bar, alternating row
          colors, and bordered cells.
  }
  \changelogentry{1.0.0}{2026-08-05}{
    \item Initial version.
    \item Added callout boxes, badge, and changelog support.
  }
\end{changelog}

\end{document}
